\section{Introduction}
\label{sec:intro}

Underwater robots are still evaluated mostly in simulation, because the
physical alternative is expensive and hard to repeat. A single wet trial with
a small remotely operated vehicle (ROV) consumes pool or vessel time, requires
an operator and a safety procedure, and yields a handful of runs whose
conditions --- turbidity, lighting, tether load, thermal drift --- are only
partly under the experimenter's control. Simulation removes those constraints
and has become the default instrument for developing and comparing underwater
perception and control: modern marine simulators provide vehicle dynamics,
camera and acoustic sensor models, and robot-middleware
interfaces~\cite{potokar2024holoocea,potokar2022holoocea}, and are routinely
used to produce the bulk of the evidence in underwater autonomy papers.

Obstacle avoidance is the task where this substitution is least safe. Whether
a vehicle passes an obstacle without contact is decided jointly by how the
obstacle is perceived, how long the perception-to-command path takes, how the
vehicle actually responds to a commanded velocity, and how well its state is
known --- and every one of those four factors is modeled approximately. A
detector that fires half a second later in reality than in simulation, or a
vehicle whose sway response is slower than modeled, changes the clearance at
which a maneuver is committed. Small, individually plausible modeling errors
therefore compose into an outcome difference exactly on the metric that
matters, namely whether the vehicle hits the obstacle.

Prior work has largely answered the transfer question and left the prediction
question open. Underwater controllers trained or tuned in simulation have been
shown to work on physical vehicles: a high-fidelity vehicle model supports
reinforcement-learning avoidance that transfers to pool
trials~\cite{zhang2025simtoreal}, a photogrammetric digital twin supports
policy validation on a physical BlueROV2 in a
harbor~\cite{mari2026deeprein}, and hardware-in-the-loop configurations let a
real vehicle computer run against a simulated
world~\cite{meyers2025testinga}. These results establish that a stack
developed in simulation \emph{can} work in water. They do not establish the
stronger property an evaluation instrument needs: that the simulator's verdict
about which method is safer, and by how much, survives contact with the
physical vehicle. A simulator can transfer a working controller and still
rank two candidate controllers in the wrong order.

Simulator predictivity is the sharper concept, and it is not new. Kadian et
al.~\cite{kadian2020simrealp} formalized it for indoor ground navigation:
run identical code in a simulator and in the corresponding real space,
correlate the per-method outcomes, and report a Sim-vs-Real Correlation
Coefficient (SRCC) that measures whether the simulator preserves the
\emph{ranking} of methods; they further showed that tuning simulator
parameters raises that correlation substantially. Truong et
al.~\cite{truong2022rethinki} then showed the converse is not automatic ---
higher-fidelity simulation does not necessarily improve sim-to-real behavior
--- which makes ``we added fidelity'' an unsupported argument until it is
measured. To our knowledge this measurement has not been carried out for
underwater obstacle avoidance, where the mismatch sources (underwater optics,
video transport latency, hydrodynamic response, absence of velocity aiding)
differ qualitatively from the indoor ground case.

This paper takes the predictivity question underwater and makes the mismatch
sources separable. We run one implementation of a camera-driven avoidance
stack --- detection, temporal obstacle estimation, local planning, body-velocity
command --- against two plants behind an identical ROS~2 interface: the
HoloOcean simulator with a native BlueROV2 Heavy vehicle, and a physical
BlueROV2 Heavy avoiding physical obstacles in a pool under independent
overhead ground truth. Rather than comparing a default simulator against
reality once, we evaluate a \emph{progressive calibration ladder}: \SZERO{}
is the uncalibrated practical model; \SONE{} matches the measured camera;
\STWO{} additionally matches the measured end-to-end timing; \STHREE{}
additionally matches the measured vehicle response and sensor behavior. Each
rung states explicitly what it calibrates and what it deliberately leaves
uncalibrated, so a change in predictivity can be attributed to a class of
mismatch instead of to fidelity in general.

The evaluation is designed to answer four questions:
\begin{itemize}
  \item[\textbf{RQ1}] How accurately does HoloOcean predict the physical
    safety and trajectory metrics of camera-driven obstacle avoidance on a
    BlueROV2?
  \item[\textbf{RQ2}] How does progressive calibration of perception, timing,
    and vehicle dynamics change that predictivity?
  \item[\textbf{RQ3}] Does simulation preserve the relative ranking of
    different avoidance strategies after physical deployment?
  \item[\textbf{RQ4}] How do simple temporal persistence and constant-velocity
    Kalman estimation differ under detector dropout, abrupt ego-motion, and
    outlier corruption?
\end{itemize}
RQ1--RQ3 are the primary questions; RQ4 is a controlled ablation of one stack
component, reported here because it is already complete and because it
supplies the method variants the ranking analysis needs.

Our provisional contributions are the following.
\begin{enumerate}
  \item A matched ROS~2 experimental architecture in which a HoloOcean
    BlueROV2 Heavy and a physical BlueROV2 Heavy are driven by the same
    perception, temporal-estimation, and local-planning code behind identical
    message interfaces, with an audited separation between runtime
    information and validation-only ground truth
    (Sec.~\ref{sec:system}). We present this as enabling methodology: paired
    pool/simulator testbeds with external ground truth already
    exist~\cite{yang2025duvindif,alineipoiana2024abluerov,meyers2025testinga}.
  \item A progressive \SZERO--\STHREE{} calibration methodology that isolates
    camera, system-timing, and vehicle-response contributions to the
    underwater simulation-to-reality gap, with each level's calibrated and
    uncalibrated quantities declared in advance
    (Sec.~\ref{sec:calibration}).
  \item A quantitative simulator-predictivity evaluation for underwater
    obstacle avoidance, combining absolute physical safety/performance error
    with preservation of method rankings across simulation and physical
    experiments, following the predictivity framing of
    Kadian et al.~\cite{kadian2020simrealp}
    (Sec.~\ref{sec:predictivity}).
  \item A controlled temporal-estimation ablation showing complementary
    behavior between simple temporal persistence and constant-velocity Kalman
    filtering: persistence predicts better through maneuver-induced
    image-motion changes, whereas the Kalman filter's gating is what confers
    outlier robustness (Sec.~\ref{sec:results}).
\end{enumerate}

We state explicitly what is \emph{not} claimed. The predictivity question,
the SRCC metric, and the idea of improving predictivity by tuning a simulator
are due to Kadian et al.~\cite{kadian2020simrealp}. HoloOcean, its ROS~2
interface, and its hardware-in-the-loop usage are due to the BYU
FRoStLab~\cite{potokar2022holoocea,potokar2024holoocea,meyers2025testinga}.
Camera-based underwater collision avoidance with filtered relative-distance
estimation on a real ROV is due to Bergantin et
al.~\cite{bergantin2024realtime}, and belief-space and
covariance-aware safety formulations are
established~\cite{han2025riskaware,han2026riskawar,li2024anobstac}. Neither
the Kalman filter nor the dynamic-window planner used here is offered as a
novel algorithm; both are deliberately standard so that the object of study
remains the simulator.
